\section{The Three Questions}

Before you shape the Weave, you must answer three questions. I have asked them of a thousand students; I have heard a thousand stammers in reply. The questions are not difficult. The answers are.

\subsection{First: What do you intend?}

The Weave does not reward vagueness. \emph{``I want to hurt him''} is a confession, not a spell. \emph{``I want to sear his sword hand with a pinpoint of flame''} is an intention. The difference is precision. A lazy intent produces a lazy effect---or, worse, an effect that is precise in ways you did not intend.

Consider the novice who declared, \emph{``I want to light the room.''} He chose [LIGHT]. He rolled. The room illuminated, yes, but so did his own hair. He spent the remainder of the examination with a faint, embarrassing glow. The Weave had answered exactly what he asked.

\subsection{Second: Which tags serve your intent?}

Tags are the verbs and nouns of the Weave's grammar. You cannot shout \emph{``fire''} and expect the universe to tremble. You must name the quality you wish to invoke.

I provide a reference in the appendix. For now, understand the families:

\begin{itemize}
    \item \textbf{Elemental.} Burning, freezing, storm, stone, wave, wind. The loud spells. The ones that leave marks on the furniture.
    \item \textbf{Force.} Strike, area, wall, bind, dispel. The sculptors of raw arcane pressure.
    \item \textbf{Mind and Veil.} Veil, scry, memory, command, fear. The spells that make the Inquisition nervous.
    \item \textbf{Life and Body.} Heal, purify, strengthen, waken, beast. The spells that keep students alive between lectures.
    \item \textbf{Space and Motion.} Leap, fold, gate, gravity. The dangerous ones. I will speak of them in Chapter Six.
    \item \textbf{Creation and Transformation.} Create, summon, transmute, animate. The ones that make the Weave sigh.
\end{itemize}

A spell may carry one tag or several. Each tag adds a layer of meaning. \emph{``I throw a [BURNING] [STRIKE] at the brigand''} is two tags. \emph{``I throw a [BURNING] [STRIKE] [AREA] at the brigand and his companions''} is three. The Weave will accommodate you. It will also hold you responsible for the extra weight.

\subsection{Third: What is the risk?}

Every casting carries a price. The price is not in silver; it is in \emph{receipts}. The Weave will answer. Whether the answer pleases you depends on the dice and the tag count.

\medskip
\noindent
\begin{tabular}{ll}
\textbf{Tags} & \textbf{The Weave's patience} \\
\midrule
One tag & The Weave is courteous. It expects little in return. \\
Two tags & The Weave is curious. It will watch you closely. \\
Three tags & The Weave is invested. It will remember. \\
Dangerous tag & The Weave is amused. It will add a surcharge. \\
\end{tabular}
\medskip

I do not speak in numbers. Numbers are for clerks. But I will tell you this: a spell with one tag is a request. A spell with four tags is a negotiation. A spell with a dangerous tag is a dare. The Weave does not lose dares.

\section{The Act of Casting}

You have answered the three questions. Now you roll.

\subsection{The Pool}

Your dice pool is the sum of your relevant Attribute and your skill in Arcana. I do not prescribe which Attribute; that depends on your approach. A scholar invoking a precise flame uses Wits. A zealot calling down lightning uses Spirit. A performer veiling their face in a crowd uses Presence. The Weave does not care which muscle you flex, only that you flex it honestly.

\subsection{The Target}

The Game Master will set a Difficulty Value. This is not a secret. A competent caster should be able to estimate it before the dice leave their palm.

\begin{itemize}
    \item A trivial working—lighting a candle, muffling a footstep—might require two successes.
    \item A standard combat spell—a blast of force, a binding of roots—requires three.
    \item A complex ritual—a short teleport, a veil over a crowd—requires four or more.
\end{itemize}

You will learn to gauge these numbers as you learn to gauge the weight of a wine jug. Experience, not theory.

\subsection{The Dice}

Roll. Count each die showing six or higher as a success. A ten counts twice. This is the Weave's opinion of your grammar.

Now look for the ones. Each die showing a one is a \emph{story beat}—a gift to the Game Master. They will spend these gifts to complicate your success or deepen your failure. Do not resent them. A spell that works perfectly is a lecture you have already memorised. A spell that works with a twist is a lesson you will remember.

\subsection{The Outcome}

Compare your successes to the Difficulty Value.

\begin{itemize}
    \item If you meet or exceed the target and rolled no ones, the Weave is satisfied. Your intent manifests cleanly.
    \item If you meet or exceed the target but rolled ones, the Weave is satisfied but curious. The Game Master will spend your story beats to add texture—a drifting spark, an echo of thunder, a moment of hesitation in an ally.
    \item If you fall short but still achieved at least one success, you have made progress. The Weave acknowledges your effort but expects more. You gain a boon—a token of momentum—to spend later.
    \item If you roll no successes at all, the Weave is not satisfied. It is not angry, merely disappointed. You gain two boons—the Weave's apology—and the Game Master will escalate the situation. Something will happen. It will not be what you wanted.
\end{itemize}

I have seen students weep at a miss. I have seen others laugh and spend their boons to re-roll a critical failure into a triumph. The dice do not judge. The Weave does not punish. The story continues.

\section{A Note on Assistance}

Two casters may weave together. One leads; the other supports. The supporting caster describes their aid—\emph{``I steady her hands''} or \emph{``I whisper the counter- rhythm to the Weave''}—and spends a boon to grant the lead caster an additional die. No more than three assistants may contribute to a single spell. The Weave has limits. So does the casting circle.

\section{The Thepyrgos Lament}

I add this for context, not instruction.

In the tower- clad streets of my home city, the Synod licenses magic. A Runekeeper's Codex is a permit; an Invoker's Symbol is a tax stamp. But a free caster—one who shapes the raw Weave without patron or writ—carries no paper trail. They are unaccountable. They are feared.

The Synod does not hunt the University don who practices Threadweaving in a sealed laboratory, surrounded by wards and graduate students. They hunt the gutter- mage who learned three tags from a dying grandmother and now improvises storms to pay his rent. The don is managed risk. The gutter- mage is wild risk. Both are free casters. The Weave does not distinguish.

If you carry this book, you may be considered the latter. I do not apologise. I only warn: cast with intention, but also with discretion. The Chain- Lanterns have long memories.

\section{Summary for the Impatient}

You intend. You name your tags. You roll Attribute plus Arcana against a Difficulty set by the Game Master. Successes are sixes and above; ones are story beats for the Game Master. Compare successes to difficulty, accept the outcome, and spend your boons.

The rest is practice. The rest is failure. The rest is reputation.

Now turn the page. The cantrips will not learn themselves.
